
\medskip
\noindent\textup{\SGARefLink{sga3:xiv:7:lemma3}{(3)}}
Suppose that \(\varphi\in E/W\) is nontrivial and that \(R\) has
one of the following types:
\[
  A_n\ (n\ge3\text{ odd}),\qquad
  D_4,\qquad D_n\ (n\ge5),\qquad E_6.
\]
These, together with the even \(A_n\) case treated separately
above, are exactly the cases in which \(E/W\) is nontrivial. The
following diagrams record the corresponding Dynkin-diagram
automorphisms.

\begin{center}
\begin{tikzpicture}[
  x=1.05cm,
  y=.92cm,
  >=Latex,
  dynkin node/.style={circle,fill=black,inner sep=2.2pt},
  dynkin edge/.style={line width=.45pt},
  dynkin arrow/.style={line width=.45pt}
]
  \node[anchor=east] at (-1.4,6.7) {\(A_n,\ n\text{ odd}\geq 3\)};
  \foreach \x/\name in {0/aone,1.3/atwo,2.6/athree,3.9/afour,5.2/afive}
    \node[dynkin node] (\name) at (\x,6.7) {};
  \draw[dynkin edge] (aone)--(atwo)--(athree)--(afour)--(afive);
  \draw[dynkin arrow,<->]
    (aone) .. controls +(1.0,1.35) and +(-1.0,1.35) .. (afive);
  \draw[dynkin arrow,<->]
    (atwo) .. controls +(.65,.72) and +(-.65,.72) .. (afour);

  \node[anchor=east] at (-1.4,4.3) {\(D_4\)};
  \node[dynkin node] (dcenter) at (2.6,4.3) {};
  \node[dynkin node] (dtop)    at (1.3,5.15) {};
  \node[dynkin node] (dbottom) at (1.3,3.45) {};
  \node[dynkin node] (dright)  at (4.0,4.3) {};
  \draw[dynkin edge] (dcenter)--(dtop);
  \draw[dynkin edge] (dcenter)--(dbottom);
  \draw[dynkin edge] (dcenter)--(dright);
  \draw[dynkin arrow,->]
    (dright) .. controls +(0,1.05) and +(.75,.48) .. (dtop);
  \draw[dynkin arrow,->]
    (dtop) .. controls +(-1.15,-.18) and +(-1.15,.18) .. (dbottom);
  \draw[dynkin arrow,->]
    (dbottom) .. controls +(.75,-.48) and +(0,-1.05) .. (dright);

  \node[anchor=east] at (-1.4,1.85) {\(D_n,\ n\geq 5\)};
  \foreach \x/\name in {0/dnone,1.25/dntwo,2.5/dnthree,3.75/dnbranch}
    \node[dynkin node] (\name) at (\x,1.85) {};
  \node[dynkin node] (dnupper) at (4.85,2.65) {};
  \node[dynkin node] (dnlower) at (4.85,1.05) {};
  \draw[dynkin edge] (dnone)--(dntwo)--(dnthree)--(dnbranch);
  \draw[dynkin edge] (dnbranch)--(dnupper);
  \draw[dynkin edge] (dnbranch)--(dnlower);
  \draw[dynkin arrow,<->]
    (dnupper) .. controls +(.8,-.3) and +(.8,.3) .. (dnlower);

  \node[anchor=east] at (-1.4,-.75) {\(E_6.\)};
  \foreach \x/\name in {0/eone,1.3/etwo,2.6/ethree,3.9/efour,5.2/efive}
    \node[dynkin node] (\name) at (\x,-.75) {};
  \node[dynkin node] (esix) at (2.6,-1.9) {};
  \draw[dynkin edge] (eone)--(etwo)--(ethree)--(efour)--(efive);
  \draw[dynkin edge] (ethree)--(esix);
  \draw[dynkin arrow,<->]
    (eone) .. controls +(1.0,1.35) and +(-1.0,1.35) .. (efive);
  \draw[dynkin arrow,<->]
    (etwo) .. controls +(.65,.72) and +(-.65,.72) .. (efour);
\end{tikzpicture}
\end{center}

Let \(S\) be a simple system of roots, not yet ordered. The group
\(E\) is the semidirect product of \(W\) with the group
\(\Psi\) of permutations of \(S\) preserving the Cartan matrix,
or equivalently the Dynkin diagram. Thus \(E/W\simeq\Psi\), and
\(\varphi\) corresponds to some \(\psi\in\Psi\).

Inspection of the diagrams shows that each \(\psi\)-orbit in
\(S\) consists of pairwise orthogonal roots. This fails for
\(A_n\) with \(n\) even, which is why that case was separated.
If \(\sigma\) is such an orbit, the reflections \(r_a\),
\(a\in\sigma\), commute; write
\[
  \rho_\sigma=\prod_{a\in\sigma}r_a.
\]
Plainly \(\rho_\sigma\) commutes with \(\psi\).

Choose a total order on \(S\) for which every \(\psi\)-orbit is a
segment, and write \(S=\{a_1,\ldots,a_n\}\). Let
\(\theta_1,\ldots,\theta_n\) and
\(c=r_1\cdots r_n\) be defined as above. The element \(c\) is a
product of the \(\rho_\sigma\), with the orbits in some order, so
it commutes with \(\psi\). Put
\[
  f=c\psi.
\]

The automorphism \(\psi\) permutes the \(\theta_i\). Indeed,
\(\psi(\theta_i)>0\), and, since \(c\) commutes with \(\psi\),
\[
  c(\psi(\theta_i))=\psi(c(\theta_i))<0.
\]
\SGARefLink{sga3:xiv:7:lemma5}{Lemma 5}\textup{\SGARefLink{sga3:xiv:7:lemma5:item:c}{(c)}} therefore gives
\(\psi(\theta_i)=\theta_j\) for some \(j\).

Finally, if \(a\in R\), \SGARefLink{sga3:xiv:7:lemma5}{Lemma 5} gives
\(a=c^m\theta_i\) for some \(m,i\). Hence
\[
  a=f^m(\psi^{-m}\theta_i)=f^m\theta_j
\]
for a suitable \(j\). Thus \(f\) and the \(\theta_i\) satisfy
both conditions of \SGARefLink{sga3:xiv:7:lemma3}{Lemma 3}, completing its proof.

\paragraph{\textbf{Remark.}}
Except in case \textup{\SGARefLink{sga3:xiv:7:lemma2}{(2)}}, one can prove that every orbit of
\(f\) has exactly
\[
  h=\frac{\operatorname{Card}(R)}{n}
\]
elements and contains exactly one \(\theta_i\). In case
\textup{(2)}, some of the \(\theta_i\) are superfluous.
